\section{引言：不是模型的问题，是配置的问题}

HumanLayer 团队在过去一年中目睹了 coding agent 以各种方式失败：忽视指令、未经提示执行危险命令、在最简单的任务上反复打转。每次失败后，本能的反应总是一样的：``我们只需要更好的模型''、``GPT-6 会修好一切''、``等训练数据里有了我用的那个库就好了''。

然而，经过数十个项目、数百次 agent 会话的实践，HumanLayer 反复得出同一个结论：\textbf{这不是模型问题（model problem），而是配置问题（configuration problem）}。

\begin{importantbox}{核心论点}
Coding agent 的质量公式：
$$\text{Coding Agent} = \text{AI Model(s)} + \text{Harness}$$
模型只是其中一半。Harness（运行时环境与配置）决定了模型能力能否被充分发挥。与其祈祷下一代模型拯救一切，不如专注于``如何从当前模型中榨取最大价值''。
\end{importantbox}

是的，模型会越来越聪明，某些现有的失败模式会消失。但正因为模型更聪明了，我们会给它们更大、更难的新问题，而它们将以意想不到的方式继续失败。\textbf{意外的失败模式是非确定性系统的根本问题}。

\begin{knowledgebox}{什么是 Harness？}
Harness 是 coding agent 的``运行时''或``外围设备''：它是模型与环境交互的所有配置点的总和。包括但不限于：skills、MCP servers、sub-agents、memory、AGENTS.md 文件等。你可以把它理解为模型的``装备系统''---模型本身是角色属性，harness 是武器、护甲和工具。
\end{knowledgebox}

\subsection{本章小结}

Coding agent 的表现不仅取决于底层模型的能力，更取决于 harness 的质量。HumanLayer 的核心主张是：当 agent 失败时，先检查 harness，再怀疑模型。


